\section{Лекция 1. Введение в C++ и первая программа}
\label{sec:lecture-01}

Первая лекция нужна не для того, чтобы выучить много синтаксиса.
Наша задача --- построить правильную картину: что такое исходный
код, кто превращает его в исполняемую программу, зачем языку типы и как уже
на первом занятии пользоваться простыми полезными абстракциями стандартной
библиотеки.

\begin{importantbox}[Что нужно уметь после лекции]
Создать файл \texttt{.cpp}, собрать его через \texttt{g++}, запустить программу,
объяснить каждую строку простого \texttt{main}, пользоваться базовыми типами,
понимать целочисленное деление и простые неявные преобразования, а также
использовать \texttt{std::string}, \texttt{std::vector} и \texttt{std::map}
на уровне их основного контракта.
\end{importantbox}

\subsection{От задачи к машинному коду}

Процессор не выполняет строку

\begin{cppcodeplain}
std::cout << "Hello!\n";
\end{cppcodeplain}

непосредственно. В конечном счёте он исполняет машинные инструкции конкретной
архитектуры. Язык программирования позволяет описывать вычисление на более
высоком уровне: через типы, операции, функции и готовые абстракции.

\begin{definition}
\textbf{Исходный код} --- текст программы, написанный программистом.
Для C++ исходные файлы часто имеют расширение \texttt{.cpp}.
\end{definition}

\begin{definition}
\textbf{Компилятор} --- программа, которая анализирует исходный код и переводит
его в другую форму. В типичном C++ toolchain результат нескольких стадий
становится машинным кодом для выбранной платформы.
\end{definition}

\begin{whybox}[Компилятор против интерпретатора]
Деление языков на «компилируемые» и «интерпретируемые» работает только как
первое приближение. Реальные реализации могут использовать байткод, JIT и
смешанные схемы. Для нас важно, что обычная C++-сборка заранее создаёт
нативный executable, а компилятор получает информацию о типах до запуска.
\end{whybox}

\subsection{Почему C++}

C++ интересен сочетанием высокоуровневых абстракций и достаточно прямого
контроля над стоимостью программы. Можно говорить «динамическая
последовательность чисел» через \texttt{std::vector}, а позже при
необходимости разбираться с размещением памяти, временем жизни и машинным
кодом.

\begin{performancebox}
Нельзя делать вывод «программа быстрая, потому что она написана на C++».
Скорость зависит от алгоритма, структуры данных, аллокаций, кэшей,
оптимизатора, ввода-вывода и железа. Полезная идея C++ --- абстракция
не обязана означать большую runtime-стоимость.
\end{performancebox}

\subsection{Практический pipeline сборки}

Одна команда

\begin{shellcode}
g++ main.cpp -o program
\end{shellcode}

скрывает несколько разных этапов. Для обучения удобно видеть цепочку

\begin{center}
\texttt{main.cpp}
$\rightarrow$
\texttt{preprocessing}
$\rightarrow$
\texttt{assembly}
$\rightarrow$
\texttt{object file}
$\rightarrow$
\texttt{linking}
$\rightarrow$
\texttt{executable}.
\end{center}

\subsubsection{Preprocessing}

До основной компиляции обрабатываются директивы препроцессора, например

\begin{cppcodeplain}
#include <iostream>
\end{cppcodeplain}

На первом уровне понимания этого достаточно: после обработки заголовка
объявления стандартного потокового ввода-вывода становятся доступны текущей
единице трансляции.

\begin{mistakebox}[Что не надо говорить]
Фраза «\texttt{\#include} подключает библиотеку» смешивает разные стадии.
Заголовок участвует в preprocessing, а связывание с уже скомпилированным кодом
происходит на стадии linking.
\end{mistakebox}

\subsubsection{Object file и linker}

После анализа C++ компилятор может сгенерировать assembly, а assembler ---
объектный файл. Объектный файл уже содержит машинный код, но часто ещё имеет
неразрешённые ссылки на другие части программы и библиотеки.

\begin{definition}
\textbf{Линкер} связывает объектные файлы и библиотеки, разрешает внешние ссылки
и создаёт итоговый executable или библиотеку.
\end{definition}

\begin{standardbox}
Стандарт C++ определяет translation units и фазы трансляции, но не требует,
чтобы реализация буквально сохраняла промежуточные \texttt{.ii},
\texttt{.s} и \texttt{.o}. Мы создаём их специально, чтобы увидеть устройство
конкретного toolchain.
\end{standardbox}

\begin{practicebox}[Эксперимент 1: разобрать сборку по частям]
Файл \texttt{examples/01\_intro\_and\_first\_program/01\_compilation\_pipeline/inspect.sh}
отдельно запускает preprocessing, генерацию assembly, создание object file и
linking. После эксперимента команда \texttt{g++} перестаёт быть магической
кнопкой.
\end{practicebox}

\subsection{Hello, World! по строкам}

\begin{cppcode}
#include <iostream>

int main() {
    std::cout << "Hello, World!\n";
    return 0;
}
\end{cppcode}

\texttt{\#include <iostream>} --- директива preprocessing.
Заголовок \texttt{iostream} предоставляет интерфейс стандартных потоков.

\texttt{main} --- особая функция обычной C++-программы: с неё начинается
выполнение пользовательского кода после необходимой работы среды выполнения.

Слово \texttt{int} означает, что функция возвращает целое значение.
Для \texttt{main} это код завершения программы.

\texttt{std::cout} --- стандартный поток вывода. Префикс \texttt{std::}
указывает пространство имён стандартной библиотеки.

\begin{warningbox}[Почему сразу пишем std::]
В старых учебных примерах часто встречается \texttt{using namespace std;}.
Конструкция допустима, но как основной стиль она скрывает происхождение имён
и может создавать конфликты. В этом конспекте стандартные имена по умолчанию
пишутся явно: \texttt{std::cout}, \texttt{std::string}, \texttt{std::vector}.
\end{warningbox}

Оператор \texttt{<<} передаёт значение в поток вывода. Позже мы увидим, что
это перегруженный оператор; пока важен его контракт.

Текст в двойных кавычках --- строковый литерал.
Последовательность \texttt{\textbackslash n} обозначает перевод строки.

\texttt{return 0;} завершает функцию и сообщает об успешном завершении.
Достижение конца \texttt{main} без явного \texttt{return} также считается
успешным, но на первых примерах оставим возврат видимым.

\subsection{Первая команда компиляции}

\begin{shellcode}
g++ -std=c++17 -Wall -Wextra -Wpedantic hello.cpp -o hello
\end{shellcode}

Здесь:

\textbf{\texttt{g++}:} C++-драйвер GCC; \textbf{\texttt{-std=c++17}:} используем стандарт C++17; \textbf{\texttt{-Wall}:} включаем большой набор полезных предупреждений; \textbf{\texttt{-Wextra}:} включаем дополнительные предупреждения; \textbf{\texttt{-Wpedantic}:} просим предупреждать об отклонениях от стандарта; \textbf{\texttt{-o hello}:} задаём имя результата.

\begin{warningbox}
Название \texttt{-Wall} не означает буквально «все предупреждения GCC».
Поэтому полезные группы warnings часто комбинируют.
\end{warningbox}

Запуск в Linux:

\begin{shellcode}
./hello
\end{shellcode}

Префикс \texttt{./} явно указывает текущий каталог. Обычно shell ищет команды
в каталогах из \texttt{PATH}, а текущий каталог туда не включён.

\subsection{Из чего состоит строка C++}

Рассмотрим

\begin{cppcodeplain}
int age = 20;
\end{cppcodeplain}

\texttt{int} --- keyword, \texttt{age} --- identifier,
\texttt{20} --- literal, \texttt{=} --- operator,
\texttt{;} --- punctuation.

Комментарии нужны для неочевидной причины, а не для пересказа кода.

\begin{cppcodeplain}
// Плохо: увеличиваем count на один.
count += 1;
\end{cppcodeplain}

Хороший комментарий объясняет, почему увеличение происходит здесь или какой
инвариант мы поддерживаем.

\subsection{Типы как информация о значениях и операциях}

\begin{cppcode}
int age = 20;
double temperature = 36.6;
char grade = 'A';
bool passed = true;
\end{cppcode}

\begin{definition}
\textbf{Тип} определяет множество допустимых значений и операций, которые
имеют смысл для объекта или выражения.
\end{definition}

Тип нужен не только человеку: он сообщает компилятору, как трактовать данные,
какие операции разрешены и какие преобразования возможны.

Для первой лекции достаточно знать группы:

\texttt{bool}; символьный \texttt{char}; целые \texttt{short}, \texttt{int}, \texttt{long}, \texttt{long long} и unsigned-варианты; \texttt{float}, \texttt{double}, \texttt{long double}.

\begin{standardbox}[Не запоминаем int = 4 байта как закон языка]
На современных desktop-системах это часто так, но стандарт не определяет
\texttt{int} фразой «четырёхбайтовое целое». Конкретный размер относится к
реализации и платформе.
\end{standardbox}

\begin{warningbox}[double не равен множеству вещественных чисел]
Floating-point хранит конечное множество представимых значений. Многие
десятичные дроби нельзя представить точно, поэтому появляются ошибки
округления. Подробно вернёмся к этому позже.
\end{warningbox}

\subsection{Инициализация и арифметика}

\begin{cppcode}
int count = 5;
double price{12.5};
\end{cppcode}

Фигурная инициализация полезна тем, что запрещает ряд сужающих преобразований,
например попытку записать \texttt{3.14} непосредственно в \texttt{int}.

Базовые арифметические операции:

\begin{cppcode}
int a = 10;
int b = 3;

int sum = a + b;       // 13
int product = a * b;   // 30
int quotient = a / b;  // 3
int rest = a % b;      // 1
\end{cppcode}

Самый важный ранний пример:

\begin{cppcode}
int x = 7 / 2;         // 3
double y = 7.0 / 2.0; // 3.5
\end{cppcode}

\begin{whybox}[Почему 7 / 2 даёт 3]
Оба операнда целые, поэтому выполняется целочисленное деление.
В \texttt{7.0 / 2.0} операнды floating-point, поэтому операция другая.
\end{whybox}

Ещё важнее:

\begin{cppcode}
int a = 1;
int b = 4;

double wrong = a / b;       // 0.0
double right = 1.0 * a / b; // 0.25
\end{cppcode}

В первой строке сначала вычисляется целочисленное \texttt{1 / 4}, получается
\texttt{0}, и только затем ноль превращается в \texttt{double}.
Тип переменной слева не меняет задним числом семантику выражения справа.

\subsection{Сравнения и простые преобразования}

Операторы \texttt{==}, \texttt{!=}, \texttt{<}, \texttt{<=}, \texttt{>},
\texttt{>=} дают значение типа \texttt{bool}.

\begin{cppcode}
int age = 20;
bool adult = age >= 18;
\end{cppcode}

C++ также выполняет неявные преобразования:

\begin{cppcode}
int whole = 5;
double real = whole;  // 5 -> 5.0
\end{cppcode}

Обратное направление может терять информацию:

\begin{cppcode}
double real = 3.9;
int whole = real;  // дробная часть теряется
\end{cppcode}

Неявное преобразование --- конкретная операция с правилами, а не «магическая
совместимость типов».

\subsection{Ввод и вывод}

\begin{cppcode}
#include <iostream>

int main() {
    int age{};

    std::cout << "Enter age: ";
    std::cin >> age;

    std::cout << "Age: " << age << '\n';
    return 0;
}
\end{cppcode}

\texttt{std::cin >> age} просит поток ввода разобрать следующее значение как
\texttt{int}. Обработку ошибочного ввода пока откладываем: для неё нужны
управляющие конструкции.

\subsection{Первая стандартная абстракция: std::string}

\begin{cppcode}
#include <string>

std::string language{"C++"};
std::cout << language << '\n';
std::cout << language.size() << '\n';
\end{cppcode}

Пока читаем \texttt{language.size()} как «спросить у объекта длину».
Полный механизм методов появится вместе с классами.

\subsection{std::vector: последовательность однотипных значений}

\begin{cppcode}
#include <vector>

std::vector<int> numbers{10, 20, 30};

std::cout << numbers[0] << '\n';
std::cout << numbers.size() << '\n';

numbers.push_back(40);
\end{cppcode}

\texttt{std::vector<int>} пока читаем буквально как
«vector, элементы которого имеют тип int». Почему тип записан в угловых
скобках, объяснит лекция о шаблонах.

\begin{importantbox}
\texttt{vector} даёт абстракцию изменяемой последовательности однотипных
элементов и сам управляет необходимыми ресурсами. Поэтому мы можем решать
задачи с последовательностями до изучения ручного управления памятью.
\end{importantbox}

Индексация начинается с нуля. \texttt{operator[]} не проверяет границы
автоматически; к последствиям вернёмся в лекции о памяти.

\subsection{std::map: ключ и значение}

\begin{cppcode}
#include <map>
#include <string>

std::map<std::string, int> frequency;

frequency["cat"] += 1;
frequency["dog"] += 1;
frequency["cat"] += 1;

std::cout << frequency["cat"] << '\n'; // 2
\end{cppcode}

\texttt{std::map<std::string, int>} читаем как соответствие
«ключ типа string $\rightarrow$ значение типа int».

\begin{whybox}[Почему новый счётчик начинается с нуля]
\texttt{map::operator[]} при отсутствии ключа создаёт для него значение.
Для \texttt{int} значение по умолчанию равно нулю, после чего
\texttt{+= 1} увеличивает счётчик.
\end{whybox}

\subsection{Не смешиваем язык, библиотеку и инструменты}

Уже сейчас существуют три разных уровня:

\textbf{Язык C++:} типы, выражения, функции, синтаксис и правила преобразований; \textbf{Стандартная библиотека:} \texttt{iostream}, \texttt{string}, \texttt{vector}, \texttt{map}; \textbf{Toolchain:} \texttt{g++}, assembler, linker, CMake, shell.

\texttt{vector} не ключевое слово языка.
Флаг \texttt{-Wall} не входит в стандарт C++.
CMake --- не компилятор.

\subsection{CMake как первая абстракция над сборкой}

Для одного файла удобно написать команду \texttt{g++}. Когда исходников и
настроек становится больше, выгоднее описать цель сборки.

\begin{cmakecode}
cmake_minimum_required(VERSION 3.16)

project(basic_abstractions LANGUAGES CXX)

add_executable(basic_abstractions main.cpp)

target_compile_features(basic_abstractions PRIVATE cxx_std_17)
\end{cmakecode}

\texttt{add\_executable} создаёт build target, а
\texttt{target\_compile\_features} задаёт требование конкретно для него.

Сборка:

\begin{shellcode}
cmake -S . -B build
cmake --build build
./build/basic_abstractions
\end{shellcode}

\texttt{-S .} задаёт каталог исходников, \texttt{-B build} --- отдельный
каталог генерируемых файлов.

\begin{practicebox}[Эксперимент 2: простые абстракции]
Программа в
\texttt{examples/01\_intro\_and\_first\_program/02\_basic\_abstractions}
использует \texttt{string}, \texttt{vector} и \texttt{map}, но специально
обходится без циклов, собственных функций, указателей и классов.
Так мы начинаем с небольшого числа уже понятных абстракций.
\end{practicebox}

\subsection{Что пока остаётся за кадром}

Мы ещё не разобрали declaration и definition, scope, условия и циклы,
категории ошибок, память и lifetime, устройство \texttt{vector},
шаблоны, классы и стоимость операций контейнеров.

Это не пробел, а карта следующих лекций.

\subsection{Проверка понимания}

После лекции нужно суметь без подсказки объяснить:

1) чем исходный файл отличается от executable; 2) зачем нужны compiler и linker; 3) что делает \texttt{\#include}; 4) почему пишется \texttt{std::cout}; 5) почему \texttt{7 / 2} не равно \texttt{3.5}; 6) что означает \texttt{std::vector<int>}; 7) что означает \texttt{std::map<std::string, int>}; 8) почему новый \texttt{map}-счётчик можно увеличить через \texttt{+= 1}; 9) чем CMake отличается от компилятора.

\subsection{Источники для сверки}

Для терминов языка и фаз трансляции используем рабочий draft C++17 N4659.
Практический GCC pipeline сверяем с \textit{An Introduction to GCC}.
Первое педагогическое объяснение compilation pipeline сопоставляем с MIT
6.096 Lecture 1. Target-based сборку сверяем с \textit{Modern CMake}.
